\documentclass[11pt]{article}
\usepackage{manual_docs_style}
\externaldocument{build/atif_transcription}[atif_transcription.pdf]

\begin{document}
\docTitle{Stage: Agentic Trajectory Evaluation}
\phantomsection\label{docstart:agentic_trajectory_evaluation}
This stage is responsible for extracting quantitative insights from one single agentic trajectory at the time in the \hyperref[docstart:atif_transcription]{ATIF trajectory format}.
By default, it computes programmatic metrics such as token usage, interaction counts, tool-call counts, file artifacts, parsed Python invocations, library usage, and output sizes.
When run with \code{--llm}, it additionally enriches each parsed Python invocation with an LLM-generated description, observable intent class, subtype, and count of meaningful numeric values printed in the matched Python output.
The programmatic evaluation is run automatically during agentic rollout.
The LLM enrichment is triggered manually in a second stage because it adds external model calls and it is slow.

The input artifact is \code{terminal-bench/runs/<run_id>/<task>/<trial>/agent_logs/atif_processed/atif_trajectory.json}.
Additionally, for traceibility and as cache, raw LLM request, response, parsing metadata, and errors are stored separately under \code{agent_logs/atif_processed/codex_description_raw/}.
If \code{codex_description_raw} already contains a valid \code{metadata.json} entry for the same parsed Python invocation and the same Codex command (model and reasoning effort included), the LLM call is skipped and the cached metadata is reused.
Note that failed entry (i.e. the one that do not match schema below) are recomputed.


The emitted value is one JSON object with name \code{trajectory_evaluation.json} with schema defined below and in the case the LLM enrichment is run

\begin{DocCode}
python workflows/trajectory/trajectory_evaluation_programmatic.py \
  [<agentic_run_id>] \
  [--llm [--llm-model MODEL_REASONING] [--parallel N]] \
  [--path PATH]
\end{DocCode}
Note that only one of \code{<agentic_run_id>} or \code{--path} may be passed.
The \code{agentic_run_id} resolves to:
\begin{DocCode}
terminal-bench/runs/<agentic_run_id>
\end{DocCode}

When the CLI is run with \code{--llm}, the configured LLM backend is called once per Python invocation and the enrichment fields are filled directly on the corresponding Python \code{detailed} entry.
The enrichment model is selected with \code{--llm-model MODEL_REASONING}.
If \code{--llm-model} is not passed, the default is \code{gpt-5.5,low}.
An additional supported model is the locally served SGLang model \code{deepseek-v4-flash}, available through an OpenAI-compatible endpoint such as \code{http://127.0.0.1:31180/v1}.


\section*{\code{trajectory_evaluation} Schema}

\newcommand{\trajectoryEvaluationNumberOfInteractionsLink}{\hyperref[sec:trajectory-evaluation-number-of-interactions]{\texttt{number\_of\_interactions}}}

\begin{LinkedDocCode}
\{
  "agentic_run_id": "<run_id>",
  "agent_id": "<agent_profile_id>",
  "atif_trajectory_path": "/abs/path/to/atif_trajectory.json",
  "tokens": \{
    "total_prompt_tokens": 1200,
    "total_completion_tokens": 250,
    "avg_prompt_tokens_per_iteration": 80.0,
    "avg_completion_tokens_per_iteration": 16.67
  \},
  "artifact_analysis": \{
    "python_scripts": \{"created": 2, "executed": 1\},
    "json_files": \{"created": 1\},
    "plots": \{"generated": 4, "inspected": 3\}
  \},
  "cost_usd": 0.123,
  "\trajectoryEvaluationNumberOfInteractionsLink": ...,
  "tool_calls": \{
    "run_shell_command": \{"count": 9\}
  \}
\}
\end{LinkedDocCode}



\section*{\code{number_of_interactions}}
\phantomsection\label{sec:trajectory-evaluation-number-of-interactions}

This is always a fixed object with the keys:
\begin{itemize}
\item \code{user}: integer, must be 1
\item \code{system}: integer, must be 0
\item \code{agent}: a dict with three keys:
\begin{DocCode}
{
  "tool_call_steps": 51,
  "tool_observation_steps": 51,
  "steps": 20
}
\end{DocCode}
\end{itemize}
Note that \code{tool_call_steps} and \code{tool_observation_steps} must have the same value, otherwise it raises an error

\begin{DocCode}
{
  "number_of_interactions": {
    "user": 1,
    "system": 0,
    "agent": {
      "tool_call_steps": 51,
      "tool_observation_steps": 51,
      "steps": 20
    }
  }
}
\end{DocCode}
In the standard case, \code{user} should be 1 and \code{system} should be 0.


\section*{\code{artifact_analysis}}
This section summarizes files created or used by the agent, including Python scripts, JSON files, and plot artifacts.
\begin{DocCode}
{
  "artifact_analysis": {
    "python_scripts": {"created": 2, "executed": 1},
    "json_files": {"created": 1},
    "plots": {
      "generated": 4,
      "inspected": 3
    }
  }
}
\end{DocCode}
A plot is counted as \code{generated} when the trajectory creates an image-like artifact, such as a \code{.png}, \code{.jpg}, \code{.jpeg}, \code{.svg}, or \code{.pdf}.
A plot is counted as \code{inspected} only when the trajectory contains a later image-capable read, open, or visualization event for that artifact.
Merely creating a plot file, listing it, or mentioning its path does not count as inspection.

\section*{\code{tool_calls}}
\code{tool_calls} is an object keyed by tools defined in \code{agents.json::<agent_entry>::available_tools}.
The structure of each object must have a minimal key called "count" and additional keys that depends on \code{agents.json::<agent_entry>::available_tools::<tool>::parse_method}.
If \code{parse_method == false}, it means that the tool input/output are not parsed and we are just counting the calls.
If \code{parse_method == "file_path"} we additionally have \code{"read"} entry which corresponds to the files that were read.
If \code{parse_method == "tree_sitter_bash"} we additionally have \code{"commands"} entry which corresponds to commands that were run.
The total sums of each value in the parse-specific key, such as \code{read} or \code{commands}, must match \code{count}.

Schematic example:
\begin{DocCode}
{
  "tool_calls": {
    "tool_name_without_parsing": {"count": 2},
    "run_shell_command": {
      "count": 9,
      "commands": {
        "ls": {"count": 2, "call_unique": 2},
        "cat": {"total_invocations": 4},
        "python": {"total_invocations": 3}
      }
    },
    "read_file": {
      "count": 4,
      "read": {
        "docs/a.tex": 3,
        "shared/schema.json": 1
      }
    }
  }
}
\end{DocCode}
Note that entries within command are further analysed as described below


\section*{commands analysis}

In the current implementation, we support \code{cat} and \code{python}.
If a command is not supported, then it will just return
\begin{DocCode}
{"count": 3, "call_unique": 2}
\end{DocCode}
Here \code{count} is the number of times the command was called, and \code{call_unique} counts the command at most once per tool call.

\subsection*{\code{cat}}

\code{cat} tracks file-creation patterns inside parsed shell snippets.

Schema:

\begin{DocCode}
{
  "cat": {
    "total_invocations": 4,
    "pattern": {
      "heredoc_write": {
        "create_files": {
          "python_files": {
            "rule.py": 1
          },
          "other": {
            "notes.txt": 2
          }
        }
      },
      "file_concat_to_stdout": 1,
      "file_concat_to_file": 1,
      "pipe_into_cat": 0,
      "other": 0
    },
    "followup_command_counts": {
      "python": 1
    },
    "execution_after_creation": {
      "executed_in_same_shell_snippet": 1,
      "not_executed_in_same_shell_snippet": 2
    },
    "heredoc_delimiters": {
      "PY": 1,
      "EOF": 1
    }
  }
}
\end{DocCode}

Notes:
\begin{itemize}
\item \code{heredoc_write} means a \code{cat} command receives heredoc input and writes it to a file, for example \code{cat << 'EOF' > script.py}.
\item \code{pattern.heredoc_write.create_files}. It can either be
\begin{itemize}
  \item \code{python_files}: All python files
  \item \code{other}: Other files
\end{itemize}
\item \code{followup_command_counts} tracks the first command after \code{cat}.
\item \code{execution_after_creation.executed_in_same_shell_snippet} increments when such a followup command exists.
\item \code{execution_after_creation.not_executed_in_same_shell_snippet} increments when no later command in the same parsed shell snippet references the created file.
\end{itemize}

\subsection*{\code{python}}

\code{python} summarizes parsed Python invocations found inside shell-parsed tools.
This block is emitted only when the observed invocation count is non-zero.

Schema with optional LLM enrichment:

\begin{DocCode}
{
  "python": {
    "total_invocations": 3,
    "failed_invocations": 1,
    "total_stdout_lines": 7,
    "total_stderr_lines": 2,
    "total_output_bytes": 142,
    "invocation_modes": {
      "dash_c": 1,
      "heredoc": 1,
      "script_path": 1,
      "other": 0
    },
    "lines_of_code": {
      "script_total": 24
    },
    "library_calls": {
      "numpy": 2,
      "pandas": 1
    },
    "detailed": {
      "dash_c_1": {
        "executions": 1,
        "lines_of_code": 1,
        "library_calls": {},
        "total_stdout_lines": 2,
        "total_stderr_lines": 0,
        "total_output_bytes": 12,
        "description": "Prints the first few time values.",
        "number_of_semantically_relevant_numbers": 3,
        "intent_class": "data_orientation",
        "specific_subtype": "data_preview"
      },
      "analyze.py": {
        "executions": 1,
        "lines_of_code": 24,
        "library_calls": {
          "numpy": 2,
          "pandas": 1
        },
        "total_stdout_lines": 5,
        "total_stderr_lines": 2,
        "total_output_bytes": 130,
        "description": "Computes summary statistics for candidate classes.",
        "number_of_semantically_relevant_numbers": 8,
        "intent_class": "numeric_inspection",
        "specific_subtype": "summary_statistics"
      }
    }
  }
}
\end{DocCode}
The LLM enrichment fields are present only when the CLI is run with \code{--llm}.

Field semantics:
\begin{itemize}
\item \code{total_invocations}: total number of parsed \code{python} command invocations.
\item \code{failed_invocations}: number of invocations whose linked observation text looks like a Python failure.
\item \code{total_stdout_lines}, \code{total_stderr_lines}, \code{total_output_bytes}: aggregate output statistics from linked observations.
\item \code{invocation_modes}: counts by invocation kind. It supports the following buckets:
\begin{itemize}
\item \code{dash_c}: commands like \code{python -c "print(1)"}.
\item \code{heredoc}: commands where Python source is passed through a heredoc.
\item \code{script_path}: commands like \code{python analyze.py}.
\item \code{other}: parsed Python invocations that do not match the previous categories.
\end{itemize}

\item \code{lines_of_code.script_total}: sum of resolved script file lengths for \code{script_path} invocations only.
\item \code{library_calls}: aggregated detected Python library usage across all analyzed code.
\item \code{detailed}: JSON object keyed by Python invocation identifier; each value summarizes one parsed invocation entry.
\end{itemize}

Each entry in \code{detailed} contains:
\begin{itemize}
\item \code{python_identifier_id} that must be one of the following
  \begin{itemize}
  \item Inline \code{python -c ...} snippets use keys like \code{dash_c_1}, \code{dash_c_2}, etc.
  \item Heredoc Python snippets use keys like \code{heredoc_1}, \code{heredoc_2}, etc.
  \item Script-path invocations use the invoked script target string, for example \code{analyze.py}.
  \item For unresolved script targets, the detailed entry still exists, but \code{lines_of_code} can be \code{null} and \code{library_calls} can be empty.
  \end{itemize}
\end{itemize}

\subsubsection*{\code{library_calls}}

By default, \code{library_calls} maps each detected library function to a total count:

\begin{DocCode}
{
  "library_calls": {
    "numpy": 2,
    "pandas": 1
  }
}
\end{DocCode}
For instance this means that two numpy functions were called and one pandas function.



\subsubsection*{LLM enrichment fields}

If the \code{codex} executable is not on \code{PATH}, set \code{TSENV_CODEX_BIN} to the Codex binary path.
The raw Codex prompt, response, error stream, and parsed metadata for each parsed Python invocation are stored next to the processed ATIF trajectory, under \code{agent_logs/atif_processed/codex_description_raw}.
This folder acts as the LLM-enrichment cache: valid entries produced by the same Codex command are skipped by default on rerun, while missing or invalid entries are recomputed.
The current schema is:
\begin{DocCode}
{
  "description": "Computes summary statistics for candidate classes.",
  "number_of_semantically_relevant_numbers": 43,
  "intent_class": "numeric_inspection",
  "specific_subtype": "summary_statistics"
}
\end{DocCode}
Where
\begin{itemize}
\item \code{description}: A short natural-language summary of what the script does.
\item \code{number_of_semantically_relevant_numbers}: The number of semantically meaningful numbers the script outputs.
\item \code{intent_class}: One of \code{data_orientation}, \code{numeric_inspection}, \code{visualization}, \code{feature_extraction}, \code{decision_construction}, \code{prediction}, \code{validation}, \code{submission}, or \code{utility}.
\item \code{specific_subtype}: A concise subtype compatible with the selected \code{intent_class}, such as \code{data_preview}, \code{summary_statistics}, \code{plot_generation}, \code{model_fit}, \code{rule_validation}, or \code{final_answer}.
\end{itemize}

In order to avoid silent errors, if the output doens't comply with this schema, we raise an error.



The number of enriched entries equals the number of parsed Python invocation entries for which LLM enrichment was run.
Each Codex request sees exactly one invocation and one matched observed output.
The resulting enrichment fields are then attached to the corresponding Python \code{detailed} entry.

The user prompt is:
\begin{DocCode}
Return only a JSON object with exactly these keys:
description, number_of_semantically_relevant_numbers, intent_class, specific_subtype
You are annotating one parsed Python invocation from an agent trajectory.
Classify the observable action performed by this invocation using only the provided
Python call and its matched observed output.
Do not infer the agent's hidden mental state or assume what happened in earlier or later trajectory steps.
Field requirements:
- description:
  In at most 20 words, describe what the Python code does and what it visibly outputs or creates.
  Avoid boilerplate such as "Runs script.py."
  Describe observable operations rather than inferred intentions.
- number_of_semantically_relevant_numbers
  Return a non-negative integer counting meaningful numeric values explicitly visible in the textual observed output.
  Count:
  - time-series measurements and time coordinates;
  - computed statistics, scores, estimates, thresholds, and physical parameters;
  - numeric predictions or class indices when they are task-relevant;
  - each visibly printed occurrence, including repeated values.
  Do not count:
  - numeric literals appearing only in the Python code;
  - numbers hidden inside files, arrays, or plots that are not textually displayed;
  - filesystem identifiers, hashes, line numbers, library versions, or warning codes;
  - wall-clock or logging timestamps unrelated to the analyzed time series;
  - traceback metadata or other incidental execution numbers.
  For plot-only or image-only output, return 0 unless numeric values are also printed as text.
  For empty output or failure-only output, return 0.
  Count only values that are actually visible; do not estimate omitted or truncated values.
- intent_class:
  Use exactly one of the following values:
  data_orientation:
    Inspects dataset structure, available files, columns, shapes, labels, class counts, or a small preview mainly to understand the data format.
  numeric_inspection:
    Makes numerical evidence visible in terminal output, including raw measurements, descriptive summaries, local time windows, event-aligned values, or comparisons.
    Use this when the main action is inspecting printed values, even if simple derived quantities such as finite-difference acceleration are computed for that inspection.
  visualization:
    Generates, saves, displays, opens, or renders a plot or image for visual analysis.
    Do not claim that a saved plot was inspected unless this invocation displays or opens it.
  feature_extraction:
    Computes reusable derived signals, detected events, change points, engineered features, or estimated physical parameters.
    Use this when producing the derived representation is the main action, rather than merely printing a few values.
  decision_construction:
    Constructs or fits a diagnostic procedure, including thresholds, hand-written heuristics, similarity methods, classifiers, regressors, or reusable rules.
  prediction:
    Applies an existing procedure to unlabeled samples and outputs candidate labels, rankings, class scores, or predictions.
  validation:
    Evaluates or refines an existing rule, model, feature, threshold, or guessed label
    using accuracy, errors, labeled examples, edge cases, or robustness checks.
  submission:
    Creates, updates, executes, or checks the required final task artifact, such as
    results.json or rule.py.
  utility:
    Performs environment setup, package checks, file management, metadata inspection,
    debugging, or code repair without itself carrying out substantive data analysis.
- specific_subtype:
  Use exactly one subtype compatible with the selected intent_class.
  For data_orientation:
    file_inventory
    schema_inspection
    shape_inspection
    data_preview
    label_inventory
  For numeric_inspection:
    raw_values
    summary_statistics
    local_time_window
    event_aligned_values
    pre_post_comparison
    cross_sample_comparison
  For visualization:
    plot_generation
    plot_display
    plot_generation_and_display
  For feature_extraction:
    signal_processing
    event_detection
    change_point_detection
    physical_parameter_estimation
    feature_engineering
  For decision_construction:
    threshold_rule
    heuristic_rule
    similarity_method
    model_fit
    rule_implementation
  For prediction:
    candidate_labeling
    batch_prediction
    ranking
  For validation:
    training_evaluation
    rule_validation
    error_analysis
    parameter_tuning
    robustness_check
    synthetic_test
  For submission:
    final_answer
    artifact_validation
    format_check
  For utility:
    debugging
    code_repair
    environment_check
    file_management
    package_check
    other_utility
Classification rules:
1. Classify the invocation's main observable action, not every operation it contains.
2. Writing results.json or rule.py is submission/final_answer, not utility.
3. Checking whether the required artifact exists, parses, imports, or executes is
   submission/artifact_validation or submission/format_check.
4. Printing labels for test samples is prediction unless the invocation writes the
   required final artifact, in which case it is submission.
5. Computing accuracy, examining mistakes, or tuning a procedure against labeled
   examples is validation.
6. Fitting a model or constructing thresholds is decision_construction, even if the
   model also prints training predictions.
7. Computing a narrow local window of values for manual inspection is
   numeric_inspection, not feature_extraction.
8. Computing derived quantities systematically across samples for later use is
   feature_extraction.
9. A plot saved without being displayed is visualization/plot_generation.
10. A plot displayed or opened is visualization/plot_display. Use
    plot_generation_and_display when both occur in the same invocation.
11. If an analytical invocation fails, classify the attempted analytical operation
    when it is clear from the code. Use utility/debugging only when the invocation's
    purpose is diagnosing or repairing the failure.
Output requirements:
- Return JSON only.
- Do not include markdown, comments, or explanatory text.
- Do not include any extra keys.
- Use the exact key names and allowed values shown above.
- number_of_semantically_relevant_numbers must be an integer.
Entry key: <entry_key>
Executions: <executions>
Calls:
- <invocation_text>
Observed output:
- <observation_text>
\end{DocCode}

A concrete example:

\begin{DocCode}
Entry key: dash_c_1
Executions: 1
Calls:
- python -c "print(df['time'].head())"
Observed output:
- 0    0.12
- 1    0.15
- 2    0.18
\end{DocCode}




\end{document}
